\documentclass[10pt,a4paper]{amsart}
\usepackage[margin=19mm]{geometry}
\usepackage{amsmath,amssymb,mathtools}
\usepackage[scr=rsfs,cal=euler]{mathalfa}
\usepackage{xfrac}
\usepackage[unicode,hidelinks,bookmarks=false]{hyperref}
\newcommand{\ie}{i.e.\ }
\newcommand{\cf}{cf.\ }
\newcommand{\resp}{respectively\ }
\newcommand{\Subsection}{\subsection}
\providecommand{\nobreakdash}{\nobreak\hbox{-}}
\setlength{\emergencystretch}{2em}
\multlinegap0pt
\usepackage{amssymb}
\usepackage[shortlabels]{enumitem}
\usepackage{bm}
\usepackage{dsfont}
\usepackage{tikz}
\usepackage{xcolor}
\newtheorem{theo}{Theorem}
\newtheorem{prop}{Proposition}
\newcommand\C{\mathbb{C}}
\providecommand{\Log}{\operatorname{Log}}
\renewcommand\Re{\operatorname{Re}}
\providecommand{\hata}{\widehat a}
\title{From the planar Ising model to quasiconformal mappings}
\author{Rémy Mahfouf$^\mathrm{A}$}
\date{Source: arXiv:2512.20361v3 (CC BY 4.0)}
\begin{document}
\maketitle

\noindent\emph{Source and licence.} From the planar Ising model to quasiconformal mappings. Version arXiv:2512.20361v3 (CC BY 4.0), distributed by arXiv under the Creative Commons Attribution 4.0 International licence, http://creativecommons.org/licenses/by/4.0/, as recorded in the arXiv licence metadata for this exact version and shown on the abstract page https://arxiv.org/abs/2512.20361v3. Full licence notice: \texttt{licenses/arxiv-2512.20361-CC-BY-4.0.txt}.\par\smallskip

\noindent\textit{Editorial scope.} Counted unit: the article's prop lem:constant-logarithm (source0.tex lines 1977--1991). The article's complete original proof of it is delivered with it (source0.tex lines 1993--2010, one \texttt{proof} environment of 18 lines). No mathematics is written, repaired or re-derived: the statement and the proof are the article's own lines, in the article's own notation, and no retained line is substituted or re-flowed.  Retained uncounted as the source-local context the proof needs: lines 855--865; lines 895--904; lines 1942--1948.  Nothing was omitted from any retained span, so the package carries no omission record.  No retained line contains a citation, so the wrapper carries no bibliography and no bibliography entry.  No retained line contains a figure environment, an image-inclusion command or drawing code, so no image is delivered and the package declares no asset dependency.  Editorial formatting only, in addition: an AMS \texttt{amsart} preamble that restates the article's own declarations and macros used by the retained lines, namely the environments \texttt{theo}, \texttt{prop} on separate counters, together with the standard packages those lines need.  The retained environments print in this document as Theo 1, Proposition 1 in order of appearance, whereas the article prints the same items with the numbering of its own full document.  The complete native source of the article as distributed in the arXiv e-print for this exact version is bundled unchanged as \texttt{sources/191558/source0.tex}.
\begin{equation}\label{eq:g_beltrami}
g_{\bar\zeta}
=
\overline{\nu}\,\overline{g_\zeta},
\qquad
\nu
:=
\frac{(\overline z_\zeta)^{1/2}}{(z_\zeta)^{1/2}}
=
\frac{\vartheta_\zeta}{z_\zeta},
\end{equation}

Beyond the small origami case, an analogous characterization of $I[f_{(a)}]$ holds when the limiting function $\vartheta \equiv 0$ is replaced by a generic $\kappa$-Lipschitz function $\vartheta$. More precisely, there exists (up to an additive constant) a unique function satisfying natural generalizations of the above properties (at least with in a regular enough class of solution to Beltrami equations). The following theorem formalizes this statement and describes the scaling limit of the two-point fermionic correlator.
\begin{theo}\label{thm:existence-generalised-log}
	Let $\vartheta$ be $\kappa<1$-Lipschitz function in $\mathbb{C}$, $\eta \in \mathbb{C}$ and $\hat{a}\in \mathbb{C}$. There exist, up to the choice of an additive constant that depends on the point $\hat{a}\in \mathbb{C}$, a unique fonction  $\zeta \mapsto g^{[\eta]}_{\vartheta,(\hat{a})}(\zeta) $ that satisfies the following properties:
	\begin{enumerate}
	\item The function \(g^{[\eta]}_{\vartheta,(\hat{a})}\) is defined on the universal cover of $\mathbb C\setminus\{\hat{a}\}$, and each of its local branches belongs to $W^{1,2}_{\mathrm{loc}}$.
		\item $g^{[\eta]}_{\vartheta,(\hat{a})}$ satisfies the conjugate Beltrami equation \eqref{eq:g_beltrami} in $\mathbb{C} \backslash \{ \hat{a} \} $.
		\item $g^{[\eta]}_{\vartheta,(\hat{a})}$ is well defined in simply connected domains of $ \mathbb{C}\backslash \{\hat{a}\}$ and has a monodromy $2i\pi \overline{\eta}$ along a positively oriented contour in $ \mathbb{C}\backslash \{\hat{a}\}$ that winds around $\hat{a}$.
		\item Each local branch of $g^{[\eta]}_{\vartheta,(\hat{a})}$ has at most logarithmic growth near $\hat{a}$ and  infinity. Moreover, $\Re[\eta \cdot g^{[\eta]}_{\vartheta,(\hat{a})}]$ produces a single-valued function on $\mathbb C\setminus\{\hat{a}\}$.	\end{enumerate}
Moreover, the function $g^{[0]}_{\vartheta,(\hat{a})}$ is constant.
	\end{theo}	

\begin{equation}\label{eq:conjugate-beltrami-local}
        g_{\bar\zeta}=q(\zeta)\,\overline{g_\zeta},
        \qquad
        q:=\overline\nu,
        \qquad
        \nu:=\frac{\vartheta_\zeta}{z_\zeta}.
\end{equation}

\begin{prop}[Constant coefficient logarithm]\label{lem:constant-logarithm}
For $\eta\in\C$, set
\begin{equation}\label{eq:constant-Aeta-Beta}
        A_\eta:=\frac{\overline\eta+q\eta}{D_q},
        \qquad
        B_\eta:=\frac{q\eta+|q|^2\overline\eta}{D_q}
        =q\,\overline{A_\eta}.
\end{equation}
Then
\begin{equation}\label{eq:constant-g-eta}
        g^{[\eta]}(\zeta)
        =A_\eta\Log W+B_\eta\overline{\Log W}
\end{equation}
solves \eqref{eq:conjugate-beltrami-local} for a constant Beltrami coefficient $q$ on $\C\setminus\{\hata\}$ and has monodromy $2\pi i\overline\eta$ around $\hata$.
\end{prop}

\begin{proof}
Given $A,B\in \mathbb{C}$, we look for a solution of the form $ g(\zeta)=A\Log W+B\overline{\Log W}$. It is straightforward to see that 
\begin{equation}\label{eq:constant-proof-derivatives}
        g_\zeta=\frac{A}{W},
        \qquad
        g_{\bar\zeta}=\frac{B}{\overline W},
        \qquad
        \overline{g_\zeta}=\frac{\overline A}{\overline W},
\end{equation}
which makes \eqref{eq:conjugate-beltrami-local} imply that $B=q\overline A$. On the other hand, along a positively oriented contour around $\hata$, $\log W$ changes by $2\pi i$ while $\overline{\log W}$ by $-2\pi i$.  Therefore the monodromy condition implies that $ A-B=\overline\eta$. Combining the two statements yields that 
\begin{equation}\label{eq:constant-proof-A}
        A=\frac{\overline\eta+q\eta}{1-|q|^2},
        \qquad
        B=q\overline A
        =\frac{q\eta+|q|^2\overline\eta}{1-|q|^2}.
\end{equation}
One can then easily check that the proposed linear combination indeed satisfies all the assumptions of Theorem \ref{thm:existence-generalised-log}. 
\end{proof}

\end{document}
